\documentclass[12pt,a4paper]{article}

\usepackage[utf8]{inputenc}
\usepackage[vietnamese]{babel}
\usepackage{amsmath,amssymb,amsthm}
\usepackage{geometry}
\usepackage{enumitem}
\usepackage[most]{tcolorbox}
\usepackage{hyperref}
\usepackage{fancyhdr}
\usepackage{tikz}
\usepackage{eso-pic}
\usepackage{xcolor}

\geometry{a4paper, margin=2cm, bottom=2.5cm}
\hypersetup{colorlinks=true, linkcolor=blue!70!black, urlcolor=blue}
\setlist[itemize]{topsep=4pt, itemsep=2pt}
\setlist[enumerate]{topsep=4pt, itemsep=2pt}

\AddToShipoutPictureFG{%
  \AtPageCenter{%
    \begin{tikzpicture}[remember picture, overlay]
      \node[rotate=45, scale=1, opacity=0.06]{\fontsize{60}{72}\selectfont\bfseries\sffamily Đặng Tấn Quí};
    \end{tikzpicture}%
  }%
}

\pagestyle{fancy}
\fancyhf{}
\fancyhead[L]{\small\textit{Chứng minh Toán học}}
\fancyhead[R]{\small\textit{Đặng Tấn Quí}}
\fancyfoot[C]{\thepage}
\fancyfoot[L]{\footnotesize\textcopyright\ Đặng Tấn Quí}
\fancyfoot[R]{\footnotesize SĐT: 0377\,122\,210}
\renewcommand{\headrulewidth}{0.4pt}
\renewcommand{\footrulewidth}{0.4pt}
\fancypagestyle{plain}{\fancyhf{}\fancyfoot[C]{\thepage}\fancyfoot[L]{\footnotesize\textcopyright\ Đặng Tấn Quí}\fancyfoot[R]{\footnotesize SĐT: 0377\,122\,210}\renewcommand{\headrulewidth}{0pt}\renewcommand{\footrulewidth}{0.4pt}}

\newtcolorbox{dinhnghia}[1][]{enhanced, breakable, colback=blue!3!white, colframe=blue!60!black, fonttitle=\bfseries, title={Định nghĩa}, #1}
\newtcolorbox{dinhly}[1][]{enhanced, breakable, colback=green!3!white, colframe=green!50!black, fonttitle=\bfseries, title={Định lý}, #1}
\newtcolorbox{tinhchat}[1][]{enhanced, breakable, colback=yellow!5!white, colframe=orange!50!black, fonttitle=\bfseries, title={Tính chất}, #1}
\newtcolorbox{phuongphap}[1][]{enhanced, breakable, colback=gray!5!white, colframe=gray!50!black, fonttitle=\bfseries, title={Phương pháp}, #1}
\newtcolorbox{luuy}[1][]{enhanced, breakable, colback=red!3!white, colframe=red!50!black, fonttitle=\bfseries, title={Lưu ý}, #1}
\newtcolorbox{congthuc}[1][]{enhanced, breakable, colback=cyan!3!white, colframe=cyan!50!black, fonttitle=\bfseries, title={Công thức}, #1}
\newtcolorbox{vidu}[1][]{enhanced, breakable, colback=violet!3!white, colframe=violet!50!black, fonttitle=\bfseries, title={Ví dụ}, #1}

\begin{document}

\thispagestyle{empty}
\begin{tikzpicture}[remember picture, overlay]
  \fill[blue!80!black] (current page.south west) rectangle (current page.north east);
  \fill[blue!60!cyan, opacity=0.8] (current page.south west) -- (current page.north west) -- ([xshift=-3cm]current page.north east) -- (current page.south east) -- cycle;
  \foreach \r/\o in {6/0.05, 4.5/0.08, 3/0.12} {\fill[white, opacity=\o] ([xshift=2cm, yshift=-2cm]current page.north east) circle (\r cm);}
  \draw[white, line width=2pt] ([yshift=-5cm]current page.north west) ++(2cm,0) -- ++(17cm,0);
  \draw[white, line width=2pt] ([yshift=-16cm]current page.north west) ++(2cm,0) -- ++(17cm,0);
  \node[anchor=west, white, font=\fontsize{28}{34}\bfseries\selectfont] at ([xshift=3cm, yshift=-7.5cm]current page.north west) {PHẦN 10: CHỨNG MINH};
  \node[anchor=west, white, font=\fontsize{24}{30}\bfseries\selectfont] at ([xshift=3cm, yshift=-9.2cm]current page.north west) {TOÁN HỌC};
  \node[anchor=west, white, font=\fontsize{16}{20}\selectfont] at ([xshift=3cm, yshift=-11cm]current page.north west) {Tài liệu luyện thi du học};
  \node[white, opacity=0.10, font=\fontsize{80}{80}\selectfont, rotate=-15] at ([xshift=-3cm, yshift=4cm]current page.center) {$\square$};
  \node[anchor=west, white, font=\Large\bfseries] at ([xshift=3cm, yshift=-13cm]current page.north west) {Biên soạn:};
  \node[anchor=west, yellow!80!white, font=\fontsize{28}{34}\bfseries\selectfont] at ([xshift=3cm, yshift=-14.5cm]current page.north west) {ĐẶNG TẤN QUÍ};
  \node[anchor=west, white!90, font=\large] at ([xshift=3cm, yshift=-18cm]current page.north west) {SĐT: 0377\,122\,210};
  \node[anchor=south, white!80, font=\small] at ([yshift=1.5cm]current page.south) {Tài liệu có bản quyền --- Cấm sao chép khi chưa được phép};
  \node[anchor=south east, white, font=\Large\bfseries] at ([xshift=-2cm, yshift=3cm]current page.south east) {\the\year};
\end{tikzpicture}
\newpage

\thispagestyle{empty}
\vspace*{5cm}
\begin{center}
  {\Large\bfseries PHẦN 10: CHỨNG MINH TOÁN HỌC}\\[8pt]
  {\large Tài liệu luyện thi du học}
\end{center}
\vspace{2cm}
\begin{tcolorbox}[enhanced, colback=red!3!white, colframe=red!60!black, fonttitle=\bfseries, title={Thông tin bản quyền}, boxrule=1.5pt]
  \textbf{Tác giả:} Đặng Tấn Quí \quad \textbf{Liên hệ:} 0377\,122\,210 \quad \textbf{Năm:} \the\year \\[6pt]
  \textbf{Bản quyền \textcopyright\ \the\year\ Đặng Tấn Quí. Bảo lưu mọi quyền.}
\end{tcolorbox}
\vfill\begin{center}\small\textit{Tài liệu lưu hành nội bộ}\end{center}
\newpage
\tableofcontents
\newpage

%% ================================================================
%%  PHẦN 10: CHỨNG MINH TOÁN HỌC
%% ================================================================
\section{Chứng minh Toán học}

%% ----------------------------------------------------------------
%%  10.1 Chứng minh quy nạp
%% ----------------------------------------------------------------
\subsection{Chứng minh quy nạp (Proof by Mathematical Induction)}

\begin{dinhnghia}[title={Nguyên lý quy nạp toán học}]
  \textbf{Nguyên lý quy nạp:} Giả sử $P(n)$ là một mệnh đề về số nguyên dương $n$. Nếu:
  \begin{enumerate}
    \item \textbf{Cơ sở quy nạp (Base Case):} $P(1)$ đúng (hoặc $P(n_0)$ đúng cho một $n_0$ nào đó).
    \item \textbf{Bước quy nạp (Inductive Step):} Giả sử $P(k)$ đúng (giả thiết quy nạp), suy ra $P(k+1)$ đúng.
  \end{enumerate}
  Thì $P(n)$ đúng với mọi $n \ge 1$ (hoặc $n \ge n_0$).
\end{dinhnghia}

\begin{phuongphap}[title={Các bước viết chứng minh quy nạp}]
  \begin{enumerate}
    \item \textbf{Phát biểu rõ:} Gọi $P(n)$ là mệnh đề cần chứng minh.
    \item \textbf{Bước cơ sở:} Chứng minh $P(1)$ (hoặc $P(n_0)$) đúng.
    \item \textbf{Bước quy nạp:} Giả sử $P(k)$ đúng cho một số nguyên $k \ge 1$ tùy ý. Chứng minh $P(k+1)$ đúng.
    \item \textbf{Kết luận:} Theo nguyên lý quy nạp, $P(n)$ đúng với mọi $n \ge 1$.
  \end{enumerate}
\end{phuongphap}

\begin{luuy}
  \begin{itemize}
    \item Trong bước quy nạp, \textbf{không được} chứng minh $P(k)$ đúng -- ta chỉ \textbf{giả thiết} nó đúng và sử dụng.
    \item Cần viết rõ ràng đâu là giả thiết quy nạp và đâu là hệ quả.
    \item \textbf{Quy nạp mạnh (Strong Induction):} Giả thiết $P(1), P(2), \ldots, P(k)$ đều đúng, suy ra $P(k+1)$.
  \end{itemize}
\end{luuy}

\begin{vidu}[title={Chứng minh tổng}]
  Chứng minh: $1 + 2 + 3 + \cdots + n = \dfrac{n(n+1)}{2}$ với mọi $n \ge 1$.

  \medskip\textbf{Chứng minh:}

  Gọi $P(n)$: $\displaystyle\sum_{k=1}^n k = \dfrac{n(n+1)}{2}$.

  \textbf{Bước cơ sở:} $P(1)$: VT $= 1$, VP $= \dfrac{1\cdot 2}{2} = 1$. Đúng. $\checkmark$

  \textbf{Bước quy nạp:} Giả sử $P(k)$ đúng: $\displaystyle\sum_{i=1}^k i = \dfrac{k(k+1)}{2}$.

  Ta cần chứng minh $P(k+1)$: $\displaystyle\sum_{i=1}^{k+1} i = \dfrac{(k+1)(k+2)}{2}$.

  \begin{align*}
    \sum_{i=1}^{k+1} i &= \underbrace{\sum_{i=1}^k i}_{\text{giả thiết}} + (k+1)
    = \frac{k(k+1)}{2} + (k+1)
    = (k+1)\!\left(\frac{k}{2} + 1\right)
    = \frac{(k+1)(k+2)}{2}. \checkmark
  \end{align*}

  Theo nguyên lý quy nạp, $P(n)$ đúng với mọi $n \ge 1$. $\blacksquare$
\end{vidu}

\begin{vidu}[title={Chứng minh bất đẳng thức}]
  Chứng minh: $2^n > n$ với mọi $n \ge 1$.

  \medskip\textbf{Chứng minh:}

  \textbf{Bước cơ sở:} $n=1$: $2^1 = 2 > 1$. Đúng. $\checkmark$

  \textbf{Bước quy nạp:} Giả sử $2^k > k$ với $k \ge 1$ tùy ý. Cần chứng minh $2^{k+1} > k+1$.
  \[
    2^{k+1} = 2 \cdot 2^k > 2k \quad (\text{vì } 2^k > k)
  \]
  Cần $2k \ge k+1$, tức $k \ge 1$. Điều này đúng theo giả thiết.
  Vậy $2^{k+1} > 2k \ge k+1$. $\checkmark$

  Kết luận: $2^n > n$ với mọi $n \ge 1$. $\blacksquare$
\end{vidu}

\begin{vidu}[title={Chứng minh chia hết}]
  Chứng minh: $n^3 + 2n$ chia hết cho $3$ với mọi $n \ge 1$.

  \medskip\textbf{Chứng minh:}

  \textbf{Bước cơ sở:} $n=1$: $1+2=3 \divisionsymbol 3$. Đúng. $\checkmark$

  \textbf{Bước quy nạp:} Giả sử $k^3+2k \divisionsymbol 3$. Ta viết $k^3+2k=3m$ với $m\in\mathbb{Z}$.

  Xét $(k+1)^3+2(k+1) = k^3+3k^2+3k+1+2k+2 = (k^3+2k) + 3k^2+3k+3 = 3m + 3(k^2+k+1)$.

  Vì $3m + 3(k^2+k+1) = 3[m+(k^2+k+1)]$, nên $(k+1)^3+2(k+1) \divisionsymbol 3$. $\checkmark$

  Kết luận: $n^3+2n \divisionsymbol 3$ với mọi $n \ge 1$. $\blacksquare$
\end{vidu}

\medskip\noindent\textbf{Các dạng bài tập quy nạp:}
\begin{enumerate}[label=\textbf{Dạng \arabic*:}, leftmargin=*, align=left]
  \item Chứng minh công thức tổng: $\displaystyle\sum_{k=1}^n f(k) = g(n)$.
  \item Chứng minh công thức tích.
  \item Chứng minh bất đẳng thức bằng quy nạp.
  \item Chứng minh tính chia hết bằng quy nạp.
  \item Quy nạp mạnh để chứng minh tính duy nhất phân tích nhân tử.
\end{enumerate}

%% ----------------------------------------------------------------
%%  10.2 Chứng minh phản chứng
%% ----------------------------------------------------------------
\subsection{Chứng minh phản chứng (Proof by Contradiction)}

\begin{phuongphap}[title={Phương pháp phản chứng (Reductio ad Absurdum)}]
  Để chứng minh mệnh đề $P$:
  \begin{enumerate}
    \item \textbf{Giả sử ngược lại:} $P$ sai (tức $\neg P$ đúng).
    \item \textbf{Suy luận:} Từ $\neg P$, dùng các bước suy luận hợp lý, dẫn đến một \textbf{mâu thuẫn} (contradiction) -- tức một phát biểu sai như $a \ne a$, hoặc mâu thuẫn với giả thiết đã biết.
    \item \textbf{Kết luận:} Do giả sử $\neg P$ dẫn đến mâu thuẫn, $\neg P$ phải sai. Vậy $P$ đúng.
  \end{enumerate}
\end{phuongphap}

\begin{luuy}
  Phản chứng đặc biệt hữu ích khi:
  \begin{itemize}
    \item Chứng minh sự tồn tại hoặc vô hạn của cái gì đó.
    \item Chứng minh số vô tỉ, số nguyên tố không hữu hạn.
    \item Chứng minh điều không thể.
  \end{itemize}
\end{luuy}

\begin{vidu}[title={$\sqrt{2}$ là số vô tỉ}]
  Chứng minh $\sqrt{2}$ là số vô tỉ.

  \medskip\textbf{Chứng minh:}

  Giả sử ngược lại: $\sqrt{2}$ hữu tỉ. Thì $\sqrt{2} = \dfrac{p}{q}$ với $p, q \in \mathbb{Z}$, $q \ne 0$, $\gcd(p,q) = 1$ (tối giản).

  $\Rightarrow 2 = \dfrac{p^2}{q^2} \Rightarrow p^2 = 2q^2$.

  Vậy $p^2$ chẵn $\Rightarrow p$ chẵn $\Rightarrow p = 2m$ với $m \in \mathbb{Z}$.

  $\Rightarrow 4m^2 = 2q^2 \Rightarrow q^2 = 2m^2 \Rightarrow q^2$ chẵn $\Rightarrow q$ chẵn.

  Nhưng cả $p$ và $q$ đều chẵn mâu thuẫn với $\gcd(p,q) = 1$.

  Vậy giả sử sai, $\sqrt{2}$ là số vô tỉ. $\blacksquare$
\end{vidu}

\begin{vidu}[title={Có vô hạn số nguyên tố}]
  Chứng minh có vô hạn số nguyên tố.

  \medskip\textbf{Chứng minh (Euclid):}

  Giả sử ngược lại: chỉ có hữu hạn số nguyên tố $p_1, p_2, \ldots, p_n$.

  Xét $N = p_1 p_2 \cdots p_n + 1$.

  Vì $N > 1$, $N$ có một ước nguyên tố $p$ nào đó. Nhưng $p \ne p_i$ với mọi $i$ (vì $N \equiv 1 \pmod{p_i}$, tức $N$ chia $p_i$ dư $1$, không thể chia hết).

  Mâu thuẫn với giả thiết $p_1, \ldots, p_n$ là tất cả các số nguyên tố.

  Vậy có vô hạn số nguyên tố. $\blacksquare$
\end{vidu}

\begin{vidu}[title={Chứng minh bất phương trình}]
  Chứng minh: Nếu $a + b = 1$ ($a, b > 0$) thì $a^2 + b^2 \ge \dfrac{1}{2}$.

  \medskip\textbf{Chứng minh (trực tiếp):}
  $a^2 + b^2 = (a+b)^2 - 2ab = 1 - 2ab$.

  Bởi AM-GM: $ab \le \dfrac{(a+b)^2}{4} = \dfrac{1}{4}$.

  Vậy $a^2 + b^2 = 1 - 2ab \ge 1 - 2 \cdot \dfrac{1}{4} = \dfrac{1}{2}$. $\blacksquare$
\end{vidu}

%% ----------------------------------------------------------------
%%  10.3 Chứng minh trực tiếp và Phép loại trừ
%% ----------------------------------------------------------------
\subsection{Chứng minh trực tiếp và Phép loại trừ (Exhaustion)}

\begin{phuongphap}[title={Chứng minh trực tiếp (Direct Proof)}]
  Để chứng minh $P \Rightarrow Q$:
  \begin{enumerate}
    \item \textbf{Giả sử $P$ đúng.}
    \item \textbf{Suy luận logic}, áp dụng định nghĩa, tính chất, định lý đã biết.
    \item \textbf{Kết luận $Q$ đúng.}
  \end{enumerate}
  Đây là phương pháp đơn giản, tự nhiên nhất.
\end{phuongphap}

\begin{phuongphap}[title={Chứng minh thuận -- nghịch chiều (Biconditional)}]
  Để chứng minh $P \Leftrightarrow Q$, cần chứng minh cả hai chiều:
  \begin{enumerate}
    \item $P \Rightarrow Q$ (chiều thuận).
    \item $Q \Rightarrow P$ (chiều nghịch).
  \end{enumerate}
\end{phuongphap}

\begin{phuongphap}[title={Chứng minh phản đảo (Proof by Contrapositive)}]
  Mệnh đề $P \Rightarrow Q$ tương đương với mệnh đề phản đảo $\neg Q \Rightarrow \neg P$.

  Đôi khi dễ hơn khi chứng minh: ``Nếu $Q$ không đúng thì $P$ không đúng.''
\end{phuongphap}

\begin{phuongphap}[title={Phép loại trừ (Proof by Exhaustion / Case Analysis)}]
  Chia tất cả các trường hợp có thể thành các \textbf{trường hợp rời rạc, không chồng chéo và bao hàm tất cả}, rồi chứng minh mệnh đề trong từng trường hợp.

  \textbf{Khi dùng:}
  \begin{itemize}
    \item Số trường hợp hữu hạn và quản lý được.
    \item Mệnh đề phụ thuộc vào chẵn/lẻ, dương/âm/bằng không, v.v.
  \end{itemize}
\end{phuongphap}

\begin{vidu}[title={Chứng minh trực tiếp -- tích hai số lẻ là số lẻ}]
  Chứng minh: Tích của hai số nguyên lẻ là số nguyên lẻ.

  \medskip\textbf{Chứng minh:}
  Gọi $a = 2m+1$ và $b = 2n+1$ (với $m,n\in\mathbb{Z}$) là hai số nguyên lẻ bất kỳ.

  $ab = (2m+1)(2n+1) = 4mn+2m+2n+1 = 2(2mn+m+n)+1$.

  Vì $2mn+m+n \in \mathbb{Z}$, $ab$ có dạng $2k+1$ với $k = 2mn+m+n$, tức $ab$ là số lẻ. $\blacksquare$
\end{vidu}

\begin{vidu}[title={Chứng minh phản đảo}]
  Chứng minh: Nếu $n^2$ lẻ thì $n$ lẻ.

  \medskip\textbf{Chứng minh (phản đảo):}

  Phản đảo: Nếu $n$ chẵn thì $n^2$ chẵn.

  Giả sử $n$ chẵn: $n = 2k$ với $k\in\mathbb{Z}$.

  $n^2 = 4k^2 = 2(2k^2)$, tức $n^2$ chẵn. Mệnh đề phản đảo đúng, vậy mệnh đề gốc đúng. $\blacksquare$
\end{vidu}

\begin{vidu}[title={Phép loại trừ -- chia trường hợp}]
  Chứng minh: $n^2 + n$ chẵn với mọi $n \in \mathbb{Z}$.

  \medskip\textbf{Chứng minh:}

  \textbf{Trường hợp 1:} $n$ chẵn, $n = 2k$: $n^2+n = 4k^2+2k = 2(2k^2+k)$ -- chẵn. $\checkmark$

  \textbf{Trường hợp 2:} $n$ lẻ, $n = 2k+1$: $n^2+n = (2k+1)^2+(2k+1) = 4k^2+4k+1+2k+1 = 4k^2+6k+2 = 2(2k^2+3k+1)$ -- chẵn. $\checkmark$

  Hai trường hợp bao phủ tất cả $n\in\mathbb{Z}$, vậy $n^2+n$ luôn chẵn. $\blacksquare$
\end{vidu}

\begin{congthuc}[title={Tóm tắt các phương pháp chứng minh}]
  \begin{tabular}{lll}
    \textbf{Phương pháp} & \textbf{Cấu trúc} & \textbf{Dùng khi nào} \\
    \hline
    Trực tiếp & $P \Rightarrow Q$ trực tiếp & Suy luận tự nhiên \\
    Phản chứng & Giả sử $\neg P$, dẫn đến mâu thuẫn & Khó chứng minh trực tiếp \\
    Phản đảo & Chứng minh $\neg Q \Rightarrow \neg P$ & Phủ định dễ hơn \\
    Quy nạp & Base + Inductive step & Mệnh đề về $n \in \mathbb{N}$ \\
    Loại trừ & Phân chia cases, chứng minh từng case & Số cases hữu hạn \\
  \end{tabular}
\end{congthuc}

\medskip\noindent\textbf{Các dạng bài tập tổng hợp:}
\begin{enumerate}[label=\textbf{Dạng \arabic*:}, leftmargin=*, align=left]
  \item Chứng minh công thức tổng, tích bằng quy nạp.
  \item Chứng minh bất đẳng thức bằng quy nạp.
  \item Chứng minh tính chia hết bằng quy nạp.
  \item Chứng minh số vô tỉ bằng phản chứng.
  \item Chứng minh mệnh đề số học bằng phản đảo.
  \item Chứng minh bằng phân chia trường hợp (loại trừ).
  \item Kết hợp nhiều phương pháp trong một bài.
\end{enumerate}

\end{document}
